\chapter{Turbulence and Turbulent Transport}\label{ch:turbulence}

Practical combustors are turbulent: the flow is an unsteady, three-dimensional tangle of eddies that enormously enhances mixing and transport. This chapter assembles the phenomenological description of turbulence needed for combustion --- the energy cascade, the integral and Kolmogorov scales, and the Reynolds number --- and then compares the turbulent time scales with the chemical time of a flame to build the Damk\"ohler and Karlovitz numbers and the premixed combustion regime diagram. It follows the turbulence portion of Lecture~8, the JMBC notes of Roekaerts, and the text of Pope~\cite{lecture8,roekaerts2010jmbc,pope2000turbulent}.

\section{The qualitative effect of turbulence}

Turbulence enhances mixing: the eddies stir reactants, products and heat together far faster than molecular diffusion could, and they wrinkle and strain a flame front, increasing its surface area and so its overall burning rate, sometimes to the point of tearing the front into disconnected pockets~\cite{lecture8,lecture4}. The extra ``turbulent diffusion'' could in principle be computed by resolving every eddy (direct numerical simulation, DNS), but for engineering flows this is unaffordable, which is why statistical modelling approaches have been developed --- the subject of Chapter~\ref{ch:turbcomb}~\cite{lecture8}. To model turbulence one first needs a language for its scales.

\section{Eddies and the energy cascade}

Turbulence is pictured as a collection of \textbf{eddies} of widely different sizes~\cite{lecture8,pope2000turbulent}. Kinetic energy is fed in at the large, energy-containing eddies of \emph{integral} length scale $\ell_t$, is transferred through an inertial range to successively smaller eddies, and is finally dissipated into heat by viscosity at the smallest, \emph{Kolmogorov} scale $\eta_k$. In statistically stationary turbulence the energy injected at the large scales equals the energy dissipated at the small scales, so the rate of dissipation $\varepsilon$ (units \si{m^2\,s^{-3}}) is set by the large eddies but is realised at the small ones~\cite{lecture8,pope2000turbulent}. This dissipation $\varepsilon$ is the same quantity that appeared as the (small) source $s_e$ in the energy equation of Chapter~\ref{ch:governing}; although negligible against the combustion heat release, it is important because it fixes the size of the smallest eddies, and hence which eddies are small enough to disturb the flame front~\cite{lecture8}.

\begin{figure}[htbp]
\centering
\includegraphics[width=0.66\linewidth]{fig_energy_cascade}
\caption{The turbulent energy cascade: energy is produced at the large (integral) scales, transferred through the inertial range, and dissipated at the Kolmogorov scale~\cite{lecture8}.}
\label{fig:fig_energy_cascade}
\end{figure}

\section{Integral scales and the turbulent Reynolds number}

The energy-containing eddies of size $\ell_t$ carry most of the kinetic energy, and their velocity scale is the r.m.s.\ fluctuation $u'\approx u'_{rms}$. The \textbf{turbulent Reynolds number} is
\begin{equation}
\Ret=\frac{u'\,\ell_t}{\nu},
\label{eq:Ret}
\end{equation}
with $\nu$ the kinematic viscosity~\cite{lecture8}. For a jet a first estimate takes $\ell_t$ of order the jet diameter and $u'$ a fraction (say \SI{10}{\percent}) of the mean injection velocity~\cite{lecture8}. The lifetime (turnover time) of a large eddy follows from the idea that during its life it moves only about its own size, giving the estimate of the dissipation rate
\begin{equation}
\varepsilon\approx\frac{u'^2}{\ell_t/u'}=\frac{u'^3}{\ell_t},
\label{eq:epsilon}
\end{equation}
the central scaling that links the large-eddy quantities to $\varepsilon$~\cite{lecture8}.

\section{The Kolmogorov scales}

The smallest eddies are characterised, following Kolmogorov's hypothesis, by the only quantities relevant there --- the dissipation $\varepsilon$ and the viscosity $\nu$ --- from which dimensional analysis builds the \textbf{Kolmogorov} length, time and velocity scales
\begin{equation}
\eta_k=\left(\frac{\nu^3}{\varepsilon}\right)^{1/4},\qquad
\tau_k=\left(\frac{\nu}{\varepsilon}\right)^{1/2},\qquad
v_k=(\nu\varepsilon)^{1/4}.
\label{eq:kolmogorov}
\end{equation}
Combining these with the integral scales and Eq.~\eqref{eq:epsilon} gives the scale-separation ratios as powers of the Reynolds number~\cite{lecture8,pope2000turbulent}:
\begin{equation}
\frac{\ell_t}{\eta_k}=\Ret^{3/4},\qquad
\frac{\tau_t}{\tau_k}=\Ret^{1/2},\qquad
\frac{v_t}{v_k}=\Ret^{1/4}.
\label{eq:scale-ratios}
\end{equation}
At $\Ret=10^4$, for example, the largest eddies are about $10^3$ times bigger and $10^2$ times longer-lived than the smallest --- the six-orders-of-magnitude span of scales anticipated in Chapter~\ref{ch:intro}~\cite{lecture8,lecture4}. Resolving all of them simultaneously is what makes DNS so costly and modelling so necessary.

\section{Turbulent transport and the closure problem}

To avoid resolving every eddy, the equations of Chapter~\ref{ch:governing} are averaged. Because the density varies strongly in a flame, combustion uses \emph{Favre} (density-weighted) averaging, $\favg{\varphi}=\avg{\rho\varphi}/\avg{\rho}$, alongside the ordinary Reynolds average $\avg{\varphi}$~\cite{lecture8,veynante2002turbulent}. Averaging the transport equations produces new unclosed correlations: the Reynolds stresses $\avg{\rho\,\ffluc{u_i}\,\ffluc{u_j}}$ in the momentum equation and the turbulent scalar fluxes $\avg{\rho\,\ffluc{u_i}\,\ffluc{Y_k}}$ in the species equations. The simplest closure assumes \emph{gradient transport} by analogy with molecular diffusion --- the turbulent flux is proportional to the mean gradient through a turbulent diffusivity $D_T$ much larger than the molecular value --- which is why turbulence is said to enhance diffusion~\cite{lecture8,roekaerts2010jmbc}. The far harder problem, the averaging of the strongly nonlinear chemical source term $\avg{\dot\omega_k}$, is the central difficulty of turbulent combustion and is deferred to Chapter~\ref{ch:turbcomb}.

\section{Comparing flame and turbulence scales: Damk\"ohler and Karlovitz numbers}

Whether and how turbulence disturbs a flame is decided by comparing time (or length) scales~\cite{transcript_l9,lecture8}. The overall chemical time of a laminar flame is the time it needs to propagate over its own thickness, $\tau_{chem}=\dL/\sLz$~\cite{lecture8}. Two dimensionless numbers then characterise the interaction:
\begin{equation}
\Da=\frac{\tau_t}{\tau_{chem}}=\frac{\ell_t\,\sLz}{u'\,\delta_L},\qquad
\Ka=\frac{\tau_{chem}}{\tau_k}=\left(\frac{u'}{\sLz}\right)^{3/2}\left(\frac{\ell_t}{\delta_L}\right)^{-1/2},
\label{eq:Da-Ka}
\end{equation}
the \textbf{Damk\"ohler number} comparing the integral (large-eddy) time to the chemical time, and the \textbf{Karlovitz number} comparing the chemical time to the smallest (Kolmogorov) eddy time~\cite{lecture8,transcript_l9}. They are not independent of the Reynolds number; in fact $\Da^2\Ka^2=(\tau_t/\tau_k)^2=\Ret^2$~\cite{lecture8}. The physical reading is intuitive: a vortex can only distort the flame, not penetrate it, if it does not act long enough or is larger than the flame thickness, so it is easier to think in length scales --- an eddy larger than the flame thickness merely wrinkles and convects the front, while only eddies smaller than the flame can enter its structure~\cite{transcript_l9}.

\section{The premixed combustion regime diagram}

Plotting the velocity ratio $u'/\sLz$ against the length ratio $\ell_t/\delta_L$ on logarithmic axes gives the \textbf{Borghi--Peters regime diagram}, divided by the lines $\Ret=1$, $\Ka=1$ and $\Da=1$ into the regimes observed experimentally (e.g.\ via OH imaging by Ciccarelli and Dorofeev)~\cite{lecture8,ciccarelli2008flame}:
\begin{description}[leftmargin=1.5em,itemsep=2pt]
\item[Wrinkled flamelets] $u'<\sLz$, $\Ka<1$, $\Da>1$: weak turbulence merely wrinkles a thin, locally laminar flame.
\item[Corrugated flamelets] $u'>\sLz$, $\Ka<1$, $\Da>1$: stronger turbulence corrugates the front and can form pockets, but the local structure is still a laminar flamelet because all eddies are larger than the flame.
\item[Thin (distributed) reaction zones] $\Ka>1$, $\Da>1$: the smallest eddies are now thinner than the preheat zone and can enter it, broadening the preheat layer, but they are still too large to disturb the very thin inner reaction zone (for methane this holds up to $\Ka\approx100$, since the inner reaction zone is roughly $\delta_r/\delta_L\approx1/10$).
\item[Well-stirred reactor] $\Ka>1$, $\Da<1$: chemistry is slower than even the large-eddy mixing, the reaction zone itself is disrupted, and the flame approaches a distributed, homogeneous reactor.
\end{description}
The two threshold lines have a clean interpretation~\cite{lecture8}: $\Da=1$ marks the point where even the large eddies are fast enough to affect the front while it passes; $\Ka=1$ where the small eddies can affect the preheat zone but the large eddies cannot; and $\Ka\approx100$ (for methane/air) where the small eddies become fast enough to perturb the inner reaction zone. In non-premixed combustion the relevant ``flame thickness'' is instead set by the strain rate and evolves as the flame develops, so the regime classification differs~\cite{lecture8}.

\begin{figure}[htbp]
\centering
\includegraphics[width=0.66\linewidth]{fig_borghi}
\caption{Borghi--Peters premixed turbulent combustion regime diagram in the $(\ell_t/\delta_L,\ u'/\sLz)$ plane~\cite{lecture8}.}
\label{fig:fig_borghi}
\end{figure}

\section{The turbulent burning velocity}

The net effect of all this wrinkling is summarised by the \textbf{turbulent burning velocity} $s_T$, the mean rate at which the turbulent flame brush advances into the reactants~\cite{lecture9}. If turbulence acts mainly to thicken the preheat zone it can be defined by analogy with the laminar case, $s_T/\sLz=\sqrt{D_T/\alpha}$; in the thin-flame limit, where the dominant effect is the increase in surface area, asymptotic analysis gives $s_T\approx\sLz(1+u'/\sLz)$ for moderate turbulence and $s_T\approx u'$ for $u'\gg\sLz$~\cite{lecture9}. In practice $s_T$ is correlated in the universal form
\begin{equation}
\frac{s_T}{\sLz}=1+b_1\left(\frac{u'}{\sLz}\right)^{b_2}\!\Le^{-b_3}\!\left(\frac{\ell_t}{\delta_L}\right)^{b_4}\!\!\cdots,
\label{eq:sT-correlation}
\end{equation}
with regime-dependent constants $b_i$ and an explicit dependence on the \emph{Lewis number of the deficient reactant} (fuel in a lean flame, oxidiser in a rich one), reflecting the differential-diffusion effects of Chapter~\ref{ch:stretch}~\cite{lecture9,ciccarelli2008flame}. Crucially, $s_T$ does not rise without bound: as $u'$ increases, $s_T$ first grows but then \emph{bends over} and eventually falls, because intense stretch begins to quench the flame locally (the \textbf{bending effect} and the quenching regime)~\cite{lecture9}. Without differential diffusion the \emph{local} flame speed is unchanged by wrinkling --- only the integrated burning rate over the larger surface changes --- and the whole flame-speed-versus-stretch framework from the laminar counterflow analysis remains valid up to turbulence intensities of roughly ten times $\sLz$ (a bound that may fail for hydrogen, with its unusually high $\sLz$)~\cite{transcript_l9}.
